\section{Commutator Subspaces in Characteristic \texorpdfstring{$p$}{p}}

Let \(A\) be a finite-dimensional associative \(k\)-algebra, where
\[
    \operatorname{char} k=p>0.
\]
Put
\[
    S(A)=\operatorname{span}_k\{\,ab-ba\mid a,b\in A\,\}.
\]
When the algebra \(A\) is clear from the context, we simply write \(S=S(A)\).

For each integer \(m\geq 0\), define
\[
    T_m(A)=\{\,r\in A\mid r^{p^m}\in S(A)\,\}.
\]
We also define
\[
    T(A)=\bigcup_{m\geq 0}T_m(A)
    =
    \{\,r\in A\mid r^{p^m}\in S(A)
    \text{ for some }m\geq 0\,\}.
\]

\begin{lemma}
    For every integer \(e\geq 0\), every \(a,b\in A\), and every
    \(\lambda,\mu\in k\), one has
    \[
        (\lambda a+\mu b)^{p^e}
        \equiv
        \lambda^{p^e}a^{p^e}+\mu^{p^e}b^{p^e}
        \pmod{S(A)}.
    \]
    More generally, for \(x_1,\ldots,x_t\in A\) and
    \(\lambda_1,\ldots,\lambda_t\in k\),
    \[
        \left(\sum_{i=1}^t\lambda_i x_i\right)^{p^e}
        \equiv
        \sum_{i=1}^t \lambda_i^{p^e}x_i^{p^e}
        \pmod{S(A)}.
    \]
\end{lemma}

\begin{proof}
    The case \(e=0\) is clear. Assume \(e\geq 1\), and put
    \[
        q=p^e.
    \]
    Let \(x_1,\ldots,x_q\in A\), and set
    \[
        w=x_1x_2\cdots x_q.
    \]
    A cyclic permutation of \(w\) differs from \(w\) by an element of \(S(A)\). Indeed,
    \[
        x_1x_2\cdots x_q-x_2\cdots x_qx_1
        =
        x_1(x_2\cdots x_q)-(x_2\cdots x_q)x_1
        \in S(A).
    \]
    Hence every word of length \(q\) is congruent modulo \(S(A)\) to each of its cyclic
    permutations.

    Now expand
    \[
        (\lambda a+\mu b)^q.
    \]
    The terms are words of length \(q\) in the two letters \(a\) and \(b\). The two
    constant words contribute
    \[
        \lambda^q a^q+\mu^q b^q.
    \]
    Every other word is a mixed word. Since \(q=p^e\), the cyclic orbit of a mixed word
    has size divisible by \(p\): its size is a divisor of \(q\), and it is not \(1\).
    All words in such an orbit are congruent modulo \(S(A)\), and their scalar
    coefficients are the same. Therefore the sum of the terms in each mixed orbit is
    congruent modulo \(S(A)\) to
    \[
        |\mathcal{O}|\,w=0,
    \]
    because \(\operatorname{char} k=p\). Thus
    \[
        (\lambda a+\mu b)^q
        \equiv
        \lambda^q a^q+\mu^q b^q
        \pmod{S(A)}.
    \]

    The formula for a finite sum follows by induction on the number of summands.
\end{proof}

\begin{lemma}
    For every fixed integer \(m\geq 0\), the subset
    \[
        T_m(A)=\{\,r\in A\mid r^{p^m}\in S(A)\,\}
    \]
    is a \(k\)-subspace of \(A\), and
    \[
        S(A)\subseteq T_m(A).
    \]
\end{lemma}

\begin{proof}
    Put
    \[
        q=p^m.
    \]

    We first prove that \(S(A)\subseteq T_m(A)\). Let \(a,b\in A\), and set
    \[
        c=ab-ba.
    \]
    By the previous lemma,
    \[
        c^q=(ab-ba)^q
        \equiv
        (ab)^q-(ba)^q
        \pmod{S(A)}.
    \]
    Here we use the fact that, in characteristic \(p\),
    \[
        (-1)^q=-1.
    \]
    Moreover,
    \[
        (ab)^q-(ba)^q
        =
        a\bigl(b(ab)^{q-1}\bigr)
        -
        \bigl(b(ab)^{q-1}\bigr)a
        \in S(A).
    \]
    Hence
    \[
        (ab-ba)^q\in S(A).
    \]

    Now let \(s\in S(A)\). Since \(S(A)\) is spanned by commutators, we may write
    \[
        s=\sum_i \lambda_i c_i,
    \]
    where each \(c_i\) is a commutator. By the previous lemma,
    \[
        s^q
        =
        \left(\sum_i\lambda_i c_i\right)^q
        \equiv
        \sum_i \lambda_i^q c_i^q
        \pmod{S(A)}.
    \]
    Each \(c_i^q\in S(A)\), so the right-hand side lies in \(S(A)\). Therefore
    \[
        s^q\in S(A).
    \]
    Hence \(s\in T_m(A)\), and so
    \[
        S(A)\subseteq T_m(A).
    \]

    It remains to prove that \(T_m(A)\) is a \(k\)-subspace. Let \(x,y\in T_m(A)\), and
    let \(\lambda,\mu\in k\). Then
    \[
        x^q,y^q\in S(A).
    \]
    Again by the previous lemma,
    \[
        (\lambda x+\mu y)^q
        \equiv
        \lambda^q x^q+\mu^q y^q
        \pmod{S(A)}.
    \]
    The right-hand side lies in \(S(A)\), and hence
    \[
        (\lambda x+\mu y)^q\in S(A).
    \]
    Therefore
    \[
        \lambda x+\mu y\in T_m(A).
    \]
    Thus \(T_m(A)\) is a \(k\)-subspace of \(A\).
\end{proof}

\begin{corollary}
    The subset
    \[
        T(A)=\{\,r\in A\mid r^{p^m}\in S(A)
        \text{ for some }m\geq 0\,\}
    \]
    is a \(k\)-subspace of \(A\), and
    \[
        S(A)\subseteq T(A).
    \]
\end{corollary}

\begin{proof}
    For every \(m\geq 0\), the previous lemma gives
    \[
        S(A)\subseteq T_m(A).
    \]
    Moreover,
    \[
        T_m(A)\subseteq T_{m+1}(A).
    \]
    Indeed, if \(x\in T_m(A)\), then \(x^{p^m}\in S(A)\). Since
    \(S(A)\subseteq T_1(A)\), we have
    \[
        \bigl(x^{p^m}\bigr)^p=x^{p^{m+1}}\in S(A).
    \]
    Hence \(x\in T_{m+1}(A)\).

    Therefore
    \[
        T(A)=\bigcup_{m\geq 0}T_m(A)
    \]
    is the union of an increasing chain of \(k\)-subspaces, and hence is itself a
    \(k\)-subspace. Also,
    \[
        S(A)\subseteq T(A).
    \]
\end{proof}

\begin{lemma}
    Let
    \[
        A=M_n(k).
    \]
    Then
    \[
        S(A)=T(A)=\{\,X\in M_n(k)\mid \operatorname{tr}(X)=0\,\}.
    \]
\end{lemma}

\begin{proof}
    Put
    \[
        V=\{\,X\in M_n(k)\mid \operatorname{tr}(X)=0\,\}.
    \]
    Since
    \[
        \operatorname{tr}(XY-YX)=0
    \]
    for all \(X,Y\in M_n(k)\), we have
    \[
        S(A)\subseteq V.
    \]

    Conversely, let \(E_{ij}\) denote the standard matrix units. If \(i\neq j\), then
    \[
        E_{ij}
        =
        E_{ii}E_{ij}-E_{ij}E_{ii}
        \in S(A).
    \]
    Also, for \(i\neq j\),
    \[
        E_{ii}-E_{jj}
        =
        E_{ij}E_{ji}-E_{ji}E_{ij}
        \in S(A).
    \]
    The trace-zero matrices are spanned by the off-diagonal matrix units \(E_{ij}\)
    with \(i\neq j\), together with the diagonal differences \(E_{ii}-E_{jj}\). Hence
    \[
        V\subseteq S(A).
    \]
    Therefore
    \[
        S(A)=V.
    \]

    By the previous corollary,
    \[
        S(A)\subseteq T(A)\subseteq A.
    \]
    Since \(S(A)=V\) has codimension \(1\) in \(A=M_n(k)\), we have either
    \[
        T(A)=S(A)
        \qquad\text{or}\qquad
        T(A)=A.
    \]
    But
    \[
        E_{11}^{p^m}=E_{11}
        \qquad\text{for every }m\geq 0,
    \]
    and
    \[
        \operatorname{tr}(E_{11})=1\neq 0.
    \]
    Hence
    \[
        E_{11}^{p^m}\notin S(A)
        \qquad\text{for every }m\geq 0.
    \]
    Therefore \(E_{11}\notin T(A)\), so \(T(A)\neq A\). It follows that
    \[
        T(A)=S(A).
    \]

    Consequently,
    \[
        S(A)=T(A)=\{\,X\in M_n(k)\mid \operatorname{tr}(X)=0\,\}.
    \]
\end{proof}

\begin{proposition}
    Let \(A\) be a finite-dimensional \(k\)-algebra, where
    \[
        \operatorname{char} k=p>0.
    \]
    Put
    \[
        S(A)=\operatorname{span}_k\{\,ab-ba\mid a,b\in A\,\},
    \]
    and
    \[
        T(A)=\{\,a\in A\mid a^{p^m}\in S(A)
        \text{ for some }m\geq 0\,\}.
    \]
    Assume that \(k\) is a splitting field for \(A\). Then
    \[
        \#\{\,\text{non-isomorphic simple }A\text{-modules}\,\}
        =
        \dim_k A-\dim_k T(A).
    \]
\end{proposition}

\begin{proof}
    Let
    \[
        J=J(A)
    \]
    be the Jacobson radical of \(A\), and set
    \[
        \overline{A}=A/J.
    \]
    Since \(A\) is finite-dimensional, \(J\) is nilpotent. Hence for every
    \(x\in J\), there exists \(m\geq 0\) such that
    \[
        x^{p^m}=0.
    \]
    Since \(0\in S(A)\), it follows that
    \[
        J\subseteq T(A).
    \]

    Let
    \[
        \pi:A\longrightarrow \overline{A}=A/J
    \]
    be the natural quotient map. We claim that
    \[
        \pi(T(A))=T(\overline{A}).
    \]

    First, since
    \[
        \pi(S(A))=S(\overline{A}),
    \]
    if \(a\in T(A)\), then for some \(m\geq 0\),
    \[
        a^{p^m}\in S(A).
    \]
    Therefore
    \[
        \pi(a)^{p^m}
        =
        \pi(a^{p^m})
        \in S(\overline{A}),
    \]
    and hence
    \[
        \pi(a)\in T(\overline{A}).
    \]
    Thus
    \[
        \pi(T(A))\subseteq T(\overline{A}).
    \]

    Conversely, suppose that
    \[
        \overline{a}=\pi(a)\in T(\overline{A}).
    \]
    Then for some \(m\geq 0\),
    \[
        \overline{a}^{p^m}\in S(\overline{A}).
    \]
    Since \(\pi(S(A))=S(\overline{A})\), there exist \(s\in S(A)\) and \(j\in J\) such that
    \[
        a^{p^m}=s+j.
    \]
    Choose \(\ell\geq 0\) sufficiently large such that
    \[
        j^{p^\ell}=0.
    \]
    By the congruence above,
    \[
        (s+j)^{p^\ell}
        \equiv
        s^{p^\ell}+j^{p^\ell}
        \pmod{S(A)}.
    \]
    Therefore
    \[
        a^{p^{m+\ell}}
        =
        (a^{p^m})^{p^\ell}
        =
        (s+j)^{p^\ell}
        \equiv
        s^{p^\ell}
        \pmod{S(A)}.
    \]
    Since \(s\in S(A)\), the lemma on \(T_\ell(A)\) gives
    \[
        s^{p^\ell}\in S(A).
    \]
    Hence
    \[
        a^{p^{m+\ell}}\in S(A).
    \]
    Therefore
    \[
        a\in T(A).
    \]
    This proves
    \[
        \pi(T(A))=T(\overline{A}).
    \]

    Since \(J\subseteq T(A)\), we get
    \[
        A/T(A)\cong \overline{A}/T(\overline{A}).
    \]
    Hence
    \[
        \dim_k A-\dim_k T(A)
        =
        \dim_k \overline{A}-\dim_k T(\overline{A}).
    \]

    Now \(k\) is a splitting field for \(A\), so the semisimple algebra
    \[
        \overline{A}=A/J(A)
    \]
    is split semisimple. Thus by the Artin--Wedderburn theorem,
    \[
        \overline{A}
        \cong
        M_{n_1}(k)\oplus\cdots\oplus M_{n_r}(k),
    \]
    where
    \[
        r=
        \#\{\,\text{non-isomorphic simple }A\text{-modules}\,\}.
    \]
    The simple \(A\)-modules and the simple \(\overline{A}\)-modules are the same,
    because \(J(A)\) annihilates every simple \(A\)-module.

    For a finite direct sum of algebras, commutators and \(p^m\)-powers are computed
    componentwise. Hence
    \[
        S(\overline{A})
        =
        S(M_{n_1}(k))\oplus\cdots\oplus S(M_{n_r}(k)).
    \]
    Also,
    \[
        T(\overline{A})
        =
        T(M_{n_1}(k))\oplus\cdots\oplus T(M_{n_r}(k)).
    \]
    Indeed, if an element belongs componentwise to the right-hand side, then each
    component satisfies the defining condition for some exponent \(p^{m_i}\); taking a
    common exponent \(p^m\) with \(m\geq \max_i m_i\) shows that the whole element
    lies in \(T(\overline{A})\).

    By the previous matrix algebra lemma,
    \[
        T(M_{n_i}(k))
        =
        \{\,X\in M_{n_i}(k)\mid \operatorname{tr}(X)=0\,\}.
    \]
    Therefore each quotient
    \[
        M_{n_i}(k)/T(M_{n_i}(k))
    \]
    is one-dimensional over \(k\). Hence
    \[
        \dim_k \overline{A}-\dim_k T(\overline{A})
        =
        \sum_{i=1}^r
        \dim_k\bigl(M_{n_i}(k)/T(M_{n_i}(k))\bigr)
        =
        r.
    \]
    Consequently,
    \[
        \dim_k A-\dim_k T(A)
        =
        r
        =
        \#\{\,\text{non-isomorphic simple }A\text{-modules}\,\}.
    \]
\end{proof}

